\clase{4}{16 de Marzo}{}

\begin{prop}
	Si $\theta \colon I \to \R$ es tal que
	\[ T(s) = (\cos \theta(s), \sin \theta(s)), \]
	entonces $k = \theta'$.
\end{prop}
\begin{proof}[Proof Other Information]
	Basta notar que:
	\begin{align*}
		k(s) &= \langle T'(s), JT(s) \rangle \\
		&= \langle (-\sin \theta(s) \cdot \theta'(s), \cos \theta(s) \cdot \theta'(s)), (-\sin \theta(s), \cos \theta(s)) \rangle \\
		&= (-\sin \theta(s))^{2} \cdot \theta'(s) + \cos \theta(s)^{2} \cdot \theta'(s) = \theta'(s)
	.\qedhere\end{align*}
\end{proof}

\begin{prop}
	$\alpha \colon I \to \R^{2}$ curva param. regular suave.
	\begin{enumerate}[(i)]
		\item Si $\phi \colon \widetilde{I} \to I$ es un difeomorfismo, entonces
		\[ k_{\alpha \circ \phi} = (\sgn \phi') \cdot k_{\alpha}. \]

		\item Si $F \colon \R^{2} \to \R^{2}$ es un mov. rígido, entonces
		\[ k_{F \circ \alpha} = (\sgn \det A) \cdot k_{\alpha}, \]
		donde $F(p) = A \cdot p + p_{0}$ y $A \colon \R^{2} \to \R^{2}$ es una transformación lineal y ortogonal.
	\end{enumerate}
\end{prop}
\begin{proof}[Proof Other Information]
	$\boxed{i.}$ Notar que
	\begin{align*}
		k_{\alpha \circ \phi} &= \frac{\langle (\alpha \circ \phi)'', J[(\alpha \circ \phi)'] \rangle}{|(\alpha \circ \phi)'|^{3}} \\
		&= \frac{\langle ((\alpha' \circ \phi) \cdot \phi')', J[(\alpha' \circ \phi) \cdot \phi']}{|(\alpha' \circ \phi) \cdot \phi'|^{3}} \\
		&= \frac{\langle (\alpha'' \circ \phi) \cdot (\phi')^{2} + (\alpha' \circ \phi) \cdot \phi'', \phi' \cdot J(\alpha' \circ \phi) \rangle}{|\alpha' \circ \phi|^{3} \cdot (\sgn \phi')^{3} \cdot (\phi')^{3}} \\
		(\alpha' \circ \phi') \perp J(\alpha' \circ \phi) &= \frac{\langle (\alpha'' \circ \phi)(\phi')^{2}, \phi' \cdot J(\alpha' \circ \phi) \rangle}{|\alpha' \circ \phi|^{3} (\phi')^{3}}
	.\end{align*}
	Por lo tanto,
	\begin{align*}
		k_{\alpha \circ \phi} &= \frac{\langle \alpha'', J \alpha' \rangle}{|\alpha'|^{3}} \cdot \phi \cdot (\sgn \phi') \\
		&= (k_{\alpha \circ \phi}) \cdot (\sgn \phi')
	.\end{align*}
	\par $\boxed{ii.}$ Similar, usando
	\begin{itemize}
		\item $(F \circ \alpha)'(t) = DF(\alpha(t)) \cdot (\alpha'(t)) = A(\alpha'(t))$.

		\item $J \circ A = (\sgn \det A) \cdot A \circ J$.
	\end{itemize}
\end{proof}

\begin{theorem}[Fundamental de las curvas planas]
	Sean $I \subset \R$ un intervalo y $k \colon I \to \R$ función suave. Entonces, existe una curva plana suave p.p.a $\alpha \colon I \to \R^{2}$ con $k_{\alpha} = k$. \par
	Además, si $\beta \colon I \to \R^{2}$ es suave, p.p.a y $k_{\beta} = k$, entonces existe mov. rígido $F \colon \R^{2} \to \R^{2}$ tal que $F \circ \beta = \alpha$.
\end{theorem}
\begin{proof}[Proof Other Information]
	\fbox{I.} Defina $\theta \colon I \to \R$ tal que 
	\[ \theta(s) = \int_{s_{0}}^{s} k(t) \ dt, \quad (s_{0} \in I) \]
	y $\alpha \colon I \to \R^{2}$ tal que
	\[ \alpha(s) = \left( \int_{s_{0}}^{s} \cos \theta(t) \ dt, \int_{s_{0}}^{s} \sin \theta(t) \ dt \right). \]
	Entonces, $\alpha$ es una curva param. suave, y
	\[ \alpha'(s) = (\cos \theta(s), \sin \theta(s)). \]
	Luego, $\alpha$ está p.p.a. Por tanto, 
	\[ k_{\alpha}(s) = \theta'(s) = k(s), \quad \forall s \in I. \]
	\par \fbox{II.} Sea $A \colon \R^{2} \to \R^{2}$ una transformación lineal ortogonal con
	\begin{itemize}
		\item $A T_{\beta}(s_{0}) = T_{\alpha}(s_{0})$,

		\item $\det A > 0$,
	\end{itemize}
	y considere
	\[ p_{0} = \alpha(s_{0}) - A \beta(s_{0}). \]
	Defina $F \colon \R^{2} \to \R^{2}$ tal que
	\[ F(p) = A p + p_{0}. \]
	Buscamos $\gamma \colon F \circ \beta = \alpha$. Tenemos:
	\begin{itemize}
		\item $k_{\gamma} = k_{\beta} = k$.

		\item $\gamma'(s) = A \beta'(s) \implies |\gamma'(s)| = 1,\ \forall s \in I$ ($A$ preserva norma).

		\item Se cumple:
		\begin{align*}
			T_{\gamma}(s_{0}) &= \gamma'(s_{0}) = A T_{\beta}(s_{0}) = T_{\alpha}(s_{0}) \\
			N_{\gamma}(s_{0}) &= JT_{\gamma}(s_{0}) = JT_{\alpha}(s_{0}) = N_{\alpha}(s_{0}) \\
			\gamma(s_{0}) &= F(\beta(s_{0})) = A \beta(s_{0}) = \alpha(s_{0}) - A \beta(s_{0}) = \alpha(s_{0})
		.\end{align*}
	\end{itemize}
	Además,
	\begin{align*}
		& \frac{d}{ds}(|T_{\gamma} - T_{\alpha}|^{2} + |N_{\gamma} - N_{\alpha}|^{2}) \\
		&= 2 \cdot \langle T_{\gamma}' - T_{\alpha}', T_{\gamma} - T_{\alpha} \rangle + 2 \cdot \langle N_{\gamma}' - N_{\alpha}', N_{\gamma} - N_{\alpha} \rangle \\
		&= 2 \cdot \langle k \cdot N_{\gamma} - k \cdot N_{\alpha}, T_{\gamma} - T_{\alpha} \rangle + 2 \cdot \langle -(k \cdot T_{\gamma} - k \cdot T_{\alpha}), N_{\gamma} - N_{\alpha} \rangle \\
		&= 0,\quad \forall s \in I
	.\end{align*}
	Como esta función se anula en $s_{0}$, obtenemos
	\[ T_{\gamma} \equiv T_{\alpha},\quad \gamma' \equiv \alpha' \quad \text{y} \quad \gamma \equiv \alpha. \qedhere \]
\end{proof}

\section{Curvas en $\R^{3}$}

Sea $\alpha \colon I \to \R^{3}$ curva suave p.p.a y denote $T(s) \coloneq \alpha'(s)$ ($T' \perp T$).

\begin{definition}[curvatura]
	La \textit{curvatura de} $\alpha$ es la función $k_{\alpha} \colon I \to \R$ tal que
	\[ k_{\alpha}(s) = |T'(s)| = |\alpha''(s)| \quad (k_{\alpha} \geq 0). \]
\end{definition}

\begin{note}
	Suponga $k_{\alpha} \neq 0$ sobre $I$. En ese caso (con $k = k_{\alpha}$), definimos:
	\[ N(s) \coloneq \frac{T'(s)}{|T'(s)|} = \frac{T'(s)}{k(s)} \]
	o sea, $T' = k \cdot N$. Además, tenemos $N' \perp N$ (pues $|N(s)| = 1,\ \forall s \in I$).
\end{note}

\begin{definition}[vector binormal]
	El \textit{vector binormal de} $\alpha$ en $s \in I$ es
	\[ B(s) = T(s) \times N(s). \]
\end{definition}
\begin{property}
	El vector binormal cumple:
	\begin{itemize}
		\item $|B| = |T| \cdot |N| \cdot \sin(\angle T,N) = 1$.

		\item $B \perp T$ y $B \perp N$.

		\item $\{T(s), N(s), B(s)\}$ es una base ortonormal positiva, que llamaremos el \textit{triedro de Frenet de} $\alpha$.

		\item $B' \perp B$.

		\item Se tiene:
		\begin{align*}
			\langle B'(s), T(s) \rangle &= \langle T'(s) \times N(s) + T(s) \times N'(s), T(s) \rangle \\
			(T(s) \times N'(s)) \perp T(s) &= \langle k(s) \cdot N(s) \times N(s), T(s) \rangle \\
			&= 0
		.\end{align*}
		Por lo tanto, $B'(s)$ y $N(s)$ son paralelos.
	\end{itemize}
\end{property}

\begin{definition}[Torsión]
	La \textit{torsión de} $\alpha$ es la función $\tau \colon I \to \R$ tal que $\tau(s) \cdot N(s)$ (o sea, $\tau(s) = \langle B'(s), N(s) \rangle$).
\end{definition}

\begin{note}
	Se tiene que
	\begin{align*}
		\langle N'(s), T(s) \rangle &= \langle N'(s), T(s) \rangle + \langle N(s), T'(s) \rangle - \langle N(s), T'(s) \rangle \\
		&= \frac{d}{ds} \rangle N(s), T(s) \rangle - k(s) \\
		&= -k(s)
	.\end{align*}
	Además, se puede comprobar que
	\[ \langle N'(s), B(s) \rangle = -\tau(s). \]
\end{note}

\begin{prop}
	Se tiene que
	\[ \begin{cases}
		T' = k \cdot N \\
		N' = -k \cdot T - \tau \cdot B \\
		B' = \tau \cdot N
	\end{cases} \]
\end{prop}
